% id: 48-18
% book: 베이직쎈 확률과 통계 · 03 확률의 뜻과 활용
% section: 유형 027 같은 것이 있는 순열을 이용하는 확률
% figure: tikz:fig-48-18
% answer: 
% answer_source: 
% variant_level: 0 (원본 전사)
% 이 파일은 build.mjs 가 items.json 에서 만든다. 고칠 때는 items.json 을 고친다.
\smash{\raisebox{1.5mm}{\dmkichul{베이직쎈 확률과 통계 48쪽 18번}}}
\noindent\begin{minipage}[t]{0.58\linewidth}\vspace{0pt}\raggedright 오른쪽 그림과 같은 도로망이 있다. $\mathrm{A}$에서 $\mathrm{B}$까지 최단 거리로 갈 때, $\mathrm{P}$를 거쳐 갈 확률을 구하시오. \dmcondfill{각 경로를 택할 확률은 같은 것으로 생각한다.}{}\end{minipage}\hfill\begin{minipage}[t]{0.4\linewidth}\vspace{0pt}\centering \resizebox{\ifdim\width>\linewidth\linewidth\else\width\fi}{!}{% 유형 027 48-18(인쇄 48쪽) — 가로 6칸 · 세로 4칸 도로망. A 는 왼쪽 위 꼭짓점, B 는 오른쪽 아래 꼭짓점,
%   P 는 A 에서 오른쪽 1칸 · 아래 2칸 지점. 스캔 화질이 낮아 TikZ 로 다시 그림.
% 점 크기: 정점 2.2pt · 파생점 2pt. 답(경우의 수·확률)은 그림에 넣지 않는다.
\begin{tikzpicture}[scale=0.42, line width=0.5pt]
  \draw[line width=0.5pt] (0,0) grid (6,4);
  \coordinate (A) at (0,4);
  \coordinate (B) at (6,0);
  \coordinate (P) at (1,2);
  \fill (A) circle (2.2pt);
  \fill (B) circle (2.2pt);
  \fill (P) circle (2pt);
  \node[left, font=\small] at (A) {$\mathrm{A}$};
  \node[right, font=\small] at (B) {$\mathrm{B}$};
  \node[above left=-2pt and -2pt, font=\small] at (P) {$\mathrm{P}$};
\end{tikzpicture}
}\end{minipage}\par
